\documentclass[a4paper,12pt]{article}

% -------------------------------------------------
% Кодировка и язык
% -------------------------------------------------
\usepackage[T2A]{fontenc}
\usepackage[utf8]{inputenc}
\usepackage[russian]{babel}

% -------------------------------------------------
% Математика
% -------------------------------------------------
\usepackage{amsmath,amssymb,mathtools}
\usepackage{icomma}

% -------------------------------------------------
% Типографика
% -------------------------------------------------
\usepackage{helvet}
\renewcommand{\familydefault}{\sfdefault}
\usepackage{microtype}
\usepackage{setspace}
\onehalfspacing
\usepackage{parskip}

% -------------------------------------------------
% Вёрстка
% -------------------------------------------------
\usepackage[
    a4paper,
    left=22mm,
    right=22mm,
    top=22mm,
    bottom=24mm
]{geometry}

\usepackage{enumitem}
\usepackage{tabularx,booktabs}
\usepackage{xcolor}
\usepackage{tikz}
\usetikzlibrary{calc}

% -------------------------------------------------
% Колонтитулы и ссылки
% -------------------------------------------------
\usepackage{fancyhdr}
\usepackage{hyperref}

\hypersetup{
    colorlinks=true,
    linkcolor=darkaccent,
    urlcolor=darkaccent
}

% -------------------------------------------------
% Цвета
% -------------------------------------------------
\definecolor{accent}{HTML}{008080}
\definecolor{darkaccent}{HTML}{004D4D}
\definecolor{textgray}{HTML}{333333}
\definecolor{lightgray}{HTML}{E8EEEE}

\color{textgray}

% -------------------------------------------------
% Заголовки и служебные команды
% -------------------------------------------------
\setlength{\parindent}{0pt}
\setlength{\headheight}{15pt}

\pagestyle{fancy}
\fancyhf{}
\fancyhead[L]{\small\textcolor{darkaccent}{Математика, 7 класс}}
\fancyhead[R]{\small\textcolor{darkaccent}{Кирилл Зиновьев}}
\fancyfoot[C]{\small\thepage}
\renewcommand{\headrulewidth}{0.4pt}

\newcommand{\sectiontitle}[1]{%
    \vspace{0.8em}
    {\Large\bfseries\textcolor{darkaccent}{#1}}\par
    \vspace{0.25em}
    \textcolor{accent}{\rule{\linewidth}{0.7pt}}
    \vspace{0.5em}
}

\newcommand{\subsectiontitle}[1]{%
    \vspace{0.6em}
    {\large\bfseries\textcolor{accent}{#1}}\par
    \vspace{0.2em}
}

\newcommand{\important}[1]{%
    \textbf{\textcolor{darkaccent}{#1}}
}

\newcommand{\answerline}{%
    \par\vspace{0.4em}
    \textcolor{lightgray}{\rule{\linewidth}{0.6pt}}
    \vspace{0.2em}
}

\begin{document}

% =================================================
% ТИТУЛЬНЫЙ ЛИСТ
% =================================================
\begin{titlepage}
\thispagestyle{empty}

\begin{center}

\vspace*{2.2cm}

{\Large\textcolor{accent}{\textbf{ЧЕК-ЛИСТ}}}

\vspace{0.8cm}

{\huge\bfseries
Бесконечные периодические дроби\\[0.25em]
и распределительное свойство умножения
}

\vspace{0.7cm}

{\large
Период десятичной дроби, деление уголком,\\
вынесение общего множителя и раскрытие скобок
}

\vfill

{\large
\textbf{Лёвин Артём Александрович}\\[0.4em]
эксклюзивно для Кирилла Зиновьева
}

\vspace{0.9cm}

{\large \today}

\vfill

\begin{tikzpicture}[x=1cm,y=1cm]
    \draw[
        accent,
        line width=1.2pt
    ]
    (0,0)
    .. controls (3.0,0.7) and (6.0,-0.7) ..
    (9,0)
    .. controls (11.5,0.55) and (13.5,0.25) ..
    (15,0.45);
\end{tikzpicture}

\vspace{0.5cm}

\end{center}
\end{titlepage}

% =================================================
% ОСНОВНОЙ ЧЕК-ЛИСТ
% =================================================

\sectiontitle{1. Бесконечные периодические десятичные дроби}

\subsectiontitle{Что называется периодом}

Бесконечная десятичная дробь называется
\important{периодической}, если начиная с некоторого места одна и та же
цифра или группа цифр повторяется бесконечно.

Например,
\[
0,3333\ldots=0,(3),
\]
\[
0,565656\ldots=0,(56),
\]
\[
0,121565656\ldots=0,121(56).
\]

Цифры в скобках называются \important{периодом дроби}.

В записи
\[
0,121(56)
\]
числа \(1,2,1\) находятся \important{до периода}, а группа \(56\)
повторяется бесконечно:
\[
0,1215656565656\ldots
\]

\subsectiontitle{Как выбрать период правильно}

В скобки записывают \important{минимальную повторяющуюся группу цифр}.

Например,
\[
0,333333\ldots=0,(3),
\]
а не
\[
0,(33).
\]

Аналогично,
\[
0,142857142857142857\ldots=0,(142857).
\]

Повторение заканчиваем фиксировать непосредственно
\important{перед началом следующего такого же блока}.

\subsectiontitle{Пример: дробь \(\frac13\)}

Выполним деление:
\[
1:3=0,333333\ldots
\]

Следовательно,
\[
\boxed{\frac13=0,(3)}.
\]

Остаток при делении снова становится равным \(1\), поэтому далее
процесс полностью повторяется.

\subsectiontitle{Пример с непериодической частью}

Рассмотрим
\[
\frac{5}{18}.
\]

Получаем
\[
\frac{5}{18}=0,27777\ldots
\]

Значит,
\[
\boxed{\frac{5}{18}=0,2(7)}.
\]

Цифра \(2\) записывается один раз, а периодом является только \(7\).

\sectiontitle{2. Как получить периодическую дробь делением уголком}

При переводе обыкновенной дроби
\[
\frac{a}{b}
\]
в десятичную выполняем деление \(a:b\).

\subsectiontitle{Алгоритм}

\begin{enumerate}[leftmargin=2em,itemsep=0.4em]
    \item Выполните обычное деление числителя на знаменатель.
    \item Если целая часть закончилась, поставьте десятичную запятую.
    \item Приписывайте к остатку нули и продолжайте деление.
    \item Следите за остатками.
    \item Как только остаток повторился, дальнейшие цифры частного тоже начнут повторяться.
    \item Выделите минимальный повторяющийся блок и запишите его в скобках.
\end{enumerate}

\subsectiontitle{Почему иногда в частном появляется ноль}

Например, при делении \(5\) на \(6\):
\[
5<6.
\]

Поэтому целая часть равна нулю:
\[
\frac56=0,\ldots
\]

После запятой рассматриваем уже \(50:6\).

Получается
\[
\frac56=0,83333\ldots=0,8(3).
\]

Если на некотором шаге текущего числа недостаточно для деления,
в соответствующем разряде частного обязательно записывается \(0\).

\sectiontitle{3. Большой период: пример с \(\frac17\)}

Одна из важных дробей:
\[
\frac17=0,142857142857142857\ldots
\]

Следовательно,
\[
\boxed{\frac17=0,(142857)}.
\]

Период содержит сразу шесть цифр:
\[
142857.
\]

Это показывает, что период не обязан состоять из одной или двух цифр.

\sectiontitle{4. Конечная дробь как бесконечная}

Конечную десятичную дробь можно продолжить нулями:
\[
1,875=1,875000000\ldots
\]

Поэтому, если условие специально требует записать число
\important{в виде бесконечной периодической десятичной дроби}, можно написать
\[
\boxed{1,875=1,875(0)}.
\]

Аналогично для целого числа:
\[
17=17,000000\ldots=17,(0).
\]

Для отрицательного числа знак сохраняется:
\[
-17=-17,(0).
\]

\important{Важно.}
Такую запись используем прежде всего тогда, когда условие явно требует
бесконечную десятичную дробь. В обычных вычислениях конечную запись
искусственно продолжать не требуется.

\sectiontitle{5. Периодическая и непериодическая дробь}

Сравните:
\[
0,12121212\ldots=0,(12)
\]
и
\[
0,101001000100001\ldots
\]

В первом числе есть повторяющийся блок \(12\), поэтому дробь периодическая.

Во втором случае длина групп нулей постоянно меняется, постоянного периода
нет. Такая дробь является бесконечной непериодической.

Для рациональных чисел действует важный факт:

\[
\boxed{
\text{Любое рациональное число имеет конечную или периодическую
десятичную запись.}
}
\]

Следовательно, бесконечная непериодическая десятичная дробь задаёт
иррациональное число.

\sectiontitle{6. Знаки при вычислениях}

При вычислениях важно сначала определить знак результата.

Например,
\[
(-0,9)\cdot(-0,1)=0,09,
\]
поскольку
\[
(-)\cdot(-)=(+).
\]

Основные правила:
\[
(+)\cdot(+)=+,\qquad
(-)\cdot(-)=+,
\]
\[
(+)\cdot(-)=-,\qquad
(-)\cdot(+)=-.
\]

Эти же правила работают при делении.

\important{Контрольный вопрос после вычисления:}
не потерян ли знак «минус»?

\sectiontitle{7. Распределительное свойство умножения}

Это одна из ключевых тем дальнейшей алгебры.

Для любых чисел или выражений \(a,b,c\):
\[
\boxed{a(b+c)=ab+ac}
\]
и
\[
\boxed{a(b-c)=ab-ac}.
\]

Внешний множитель умножается
\important{на каждое слагаемое внутри скобок}.

Например,
\[
0,3(7a-3b)
=
0,3\cdot7a-0,3\cdot3b.
\]

Выполняем умножение коэффициентов:
\[
\boxed{0,3(7a-3b)=2,1a-0,9b}.
\]

\subsectiontitle{Типичная ошибка}

Нельзя писать
\[
a(a-b)=a^2-b.
\]

Правильно:
\[
\boxed{a(a-b)=a^2-ab}.
\]

Множитель \(a\) обязан умножиться и на \(a\), и на \(b\).

\sectiontitle{8. Вынесение общего множителя за скобки}

Распределительное свойство можно применять в обратную сторону:
\[
ab+ac=a(b+c),
\]
\[
ab-ac=a(b-c).
\]

Это называется \important{вынесением общего множителя за скобки}.

Например,
\[
47\cdot67-47\cdot33.
\]

Общий множитель --- \(47\):
\[
47\cdot67-47\cdot33
=
47(67-33).
\]

Получаем
\[
47(67-33)=47\cdot34.
\]

Если требуется доказать делимость на \(17\), уже хорошо видно:
\[
34=2\cdot17.
\]

Следовательно,
\[
47\cdot34
=
47\cdot2\cdot17,
\]
а значит, исходное выражение делится на \(17\).

\sectiontitle{9. Зачем выносить общий множитель}

Главная цель этого приёма --- упростить выражение
\important{до выполнения громоздких вычислений}.

Например,
\[
53\cdot69-53\cdot31.
\]

Не нужно отдельно вычислять два больших произведения.

Выносим \(53\):
\[
53(69-31)=53\cdot38.
\]

Если далее требуется проверить делимость на \(19\), замечаем:
\[
38=2\cdot19.
\]

Тогда
\[
53\cdot38=53\cdot2\cdot19,
\]
поэтому число делится на \(19\).

\sectiontitle{10. Корректное оформление доказательства делимости}

Если в исходном выражении деления на некоторое число не было,
нельзя внезапно приписывать знак деления внутри цепочки равенств.

Корректно:
\[
53\cdot69-53\cdot31
=
53(69-31)
=
53\cdot38
=
53\cdot2\cdot19.
\]

Отсюда делаем отдельный вывод:

\[
\boxed{
53\cdot69-53\cdot31
\text{ делится на }19.
}
\]

Или отдельно проверяем частное:
\[
\frac{53\cdot38}{19}=106.
\]

\sectiontitle{11. Мини-чек-лист перед сдачей работы}

\begin{enumerate}[leftmargin=2em,itemsep=0.45em]
    \item Верно ли определён период десятичной дроби?
    \item В скобки взята минимальная повторяющаяся группа?
    \item Не включена ли в период цифра, которая относится только к непериодической части?
    \item При делении уголком записаны все необходимые нули в частном?
    \item Сохранён ли знак отрицательного числа?
    \item При раскрытии скобок внешний множитель умножен на каждый член?
    \item При вынесении общего множителя проверено содержимое скобок обратным раскрытием?
    \item Можно ли сначала упростить выражение и избежать больших вычислений?
    \item В цепочке равенств не появились ли действия или числа «из воздуха»?
\end{enumerate}

% =================================================
% ТРЕНИРОВОЧНЫЙ БЛОК
% =================================================
\clearpage
\sectiontitle{Тренировочный блок}

\textbf{Путь А.}
Задания расположены по нарастающей сложности.
Решения оформляйте последовательно, особенно при делении уголком
и работе со скобками.

\subsectiontitle{Средний уровень}

\textbf{1.}
Представьте дробь в виде бесконечной периодической десятичной дроби:
\[
\frac56.
\]
\answerline

\textbf{2.}
Представьте дробь в виде бесконечной периодической десятичной дроби:
\[
\frac7{12}.
\]
\answerline

\textbf{3.}
Запишите число
\[
2,375
\]
в виде бесконечной периодической десятичной дроби с периодом \(0\).
\answerline

\subsectiontitle{Уровень выше среднего}

\textbf{4.}
Представьте дробь
\[
-\frac{13}{30}
\]
в виде бесконечной периодической десятичной дроби.
\answerline

\textbf{5.}
Представьте дробь
\[
\frac5{27}
\]
в виде бесконечной периодической десятичной дроби.
\answerline

\textbf{6.}
Для дроби
\[
\frac{13}{60}
\]
найдите десятичную запись, укажите непериодическую часть и период.
\answerline

\textbf{7.}
Не выполняя прямого умножения больших чисел, докажите, что
\[
73\cdot41-39\cdot41
\]
делится на \(17\). Найдите частное от деления полученного числа на \(17\).
\answerline

\textbf{8.}
Раскройте скобки и приведите подобные множители:
\[
0,4(9a-5b).
\]
\answerline

\textbf{9.}
Вынесите наибольший общий числовой множитель за скобки:
\[
28x-42y.
\]
\answerline

\subsectiontitle{Сложное задание}

\textbf{10.}
Упростите выражение
\[
0,6(15a-10b)-3(2a-b).
\]

Затем докажите, что при целых \(a\) и \(b\) значение выражения
делится на \(3\).
\answerline

% =================================================
% ОТВЕТЫ
% =================================================
\clearpage
\sectiontitle{Ответы к тренировочному блоку}

\renewcommand{\arraystretch}{1.45}

\begin{table}[ht]
\centering
\caption{Краткие ответы}
\begin{tabularx}{\linewidth}{@{}cX@{}}
\toprule
\textbf{№} & \textbf{Ответ} \\
\midrule

1 &
\[
\frac56=0,8(3).
\]
\\

2 &
\[
\frac7{12}=0,58(3).
\]
\\

3 &
\[
2,375=2,375(0).
\]
\\

4 &
\[
-\frac{13}{30}=-0,4(3).
\]
\\

5 &
\[
\frac5{27}=0,(185).
\]
\\

6 &
\[
\frac{13}{60}=0,21(6).
\]
Непериодическая часть после запятой: \(21\);
период: \(6\).
\\

7 &
\[
73\cdot41-39\cdot41
=
41(73-39)
=
41\cdot34
=
41\cdot2\cdot17.
\]
Делится на \(17\); частное:
\[
82.
\]
\\

8 &
\[
0,4(9a-5b)=3,6a-2b.
\]
\\

9 &
\[
28x-42y=14(2x-3y).
\]
\\

10 &
\[
0,6(15a-10b)-3(2a-b)
=
9a-6b-6a+3b
=
3a-3b
=
3(a-b).
\]
Следовательно, при целых \(a,b\) значение выражения кратно \(3\).
\\

\bottomrule
\end{tabularx}
\end{table}

\vfill

\begin{center}
\textcolor{accent}{\rule{0.65\linewidth}{0.7pt}}

\vspace{0.5em}

{\small
Лёвин Артём Александрович
}
\end{center}

\end{document}